\documentclass[10pt,a4paper]{amsart}
\usepackage[margin=19mm]{geometry}
\usepackage{amsmath,amssymb,mathtools}
\usepackage[scr=rsfs,cal=euler]{mathalfa}
\usepackage{xfrac}
\usepackage[unicode,hidelinks,bookmarks=false]{hyperref}
\newcommand{\ie}{i.e.\ }
\newcommand{\cf}{cf.\ }
\newcommand{\resp}{respectively\ }
\newcommand{\Subsection}{\subsection}
\providecommand{\nobreakdash}{\nobreak\hbox{-}}
\setlength{\emergencystretch}{2em}
\multlinegap0pt
\usepackage{bm}
\usepackage{bbm}
\usepackage{dsfont}
\usepackage{xspace}
\usepackage{cancel}
\usepackage{mathtools}
\usepackage{amsthm}
\makeatletter
\@ifundefined{theorem}{\newtheorem{theorem}{Theorem}}{}
\@ifundefined{corollary}{\newtheorem{corollary}[theorem]{Corollary}}{}
\makeatother
\title[Beyond Murray's Law: Non-Universal Branching Exponents from ]{Beyond Murray's Law: Non-Universal Branching Exponents from Vessel-Wall Metabolic Costs}
\author{Riccardo Marchesi}
\date{Source: arXiv:2603.13687v3 (CC BY 4.0)}
\begin{document}
\maketitle

\noindent\emph{Source and licence.} Beyond Murray's Law: Non-Universal Branching Exponents from Vessel-Wall Metabolic Costs. Version arXiv:2603.13687v3 (CC BY 4.0), distributed under the Creative Commons Attribution 4.0 International licence, as recorded in the arXiv licence metadata for this exact version and shown on the abstract page https://arxiv.org/abs/2603.13687v3. Full rights notice: licenses/arxiv-2603.13687-CC-BY-4.0.txt.\par\smallskip

\noindent\textit{Editorial scope.} Counted unit: the Corollary whose source label is \texttt{cor:asym} in the article (source0.tex lines 505--512, source label \texttt{cor:asym}), delivered together with the article's own complete original proof of it (source0.tex lines 514--574, one \texttt{proof} environment of 61 lines). No mathematics is written, repaired or re-derived: statement and proof are the article's own text line for line. Required context retained uncounted: eq:phi\_compact (lines 252--255); thm:unique (lines 334--338); cor:murray\_uniqueness (lines 576--588). Every definition, display and label of those lines is retained unchanged. Omitted from this package and not redistributed: 1--251, 256--333, 339--504, 513--513, 575--575, 589--1452. Nothing inside a retained excerpt is omitted or altered: every retained body is delivered exactly as the article's lines are written, so no omission receipt of any kind accompanies this package. Citations: the retained excerpts carry no citation command at all, so no bibliography entry of the article is transcribed and no reference is redistributed. Reference adaptation: none was needed and none was made; every reference inside the delivered unit resolves inside it, so no unresolved reference or citation is produced. Printed slips kept literal: no printed word, symbol, hypothesis or number of the counted statement or of its proof was altered. Layout and class: the article's own class is \texttt{article} with packages including geometry, amsmath, amssymb, amsthm, graphicx, booktabs, hyperref, authblk, mathpazo, microtype, float; the wrapper is the lane's pinned \texttt{amsart} editorial wrapper at 10pt on A4, which restates the theorem-like environments the retained text needs and adds the article's own macro definitions that the retained text needs (none), with extra wrapper packages (bm, bbm, dsfont, xspace, cancel, mathtools). The delivered numbering is the unit's own. Assets: no retained span contains a figure environment, a graphics-inclusion command, a PDF-page inclusion, a table or a TikZ picture; no image file is copied into this package, the package declares no asset dependency and no image file is redistributed with it. Licence: version arXiv:2603.13687v3 of this article is distributed under the Creative Commons Attribution 4.0 International licence, as recorded in the arXiv licence metadata for this exact version and shown on the abstract page https://arxiv.org/abs/2603.13687v3. A full rights notice is delivered at licenses/arxiv-2603.13687-CC-BY-4.0.txt. Characters: every retained excerpt is pure ASCII, so the delivered engine, XeLaTeX with its default Latin Modern text font, reports no missing character.
\begin{equation}
  \Phi(r,Q) = A(Q)\,r^{-4} + B\,r^2 + C\,r^{1+p}\,,
  \label{eq:phi_compact}
\end{equation}

\begin{theorem}[Existence and uniqueness of the optimal radius]
\label{thm:unique}
For every $Q>0$ and every $A,B,C,p>0$, the function $r\mapsto\Phi(r,Q)$
has a unique global minimum $r^*(Q)$ on $(0,+\infty)$.
\end{theorem}

\begin{corollary}[Uniqueness of Murray scaling]
\label{cor:murray_uniqueness}
Among all cost functions of the form~\eqref{eq:phi_compact} with $B,C\geq0$ (not
both zero), Murray's cubic branching law ($\alpha=3$) is the unique member
admitting a universal, scale-independent exponent. Any deviation introduced by
the biological wall cost ($p<1$) structurally breaks the Euler homogeneity of
the cost function (rendering it inhomogeneous with incommensurate scaling),
rendering the branching exponent dependent on the absolute flow scale $Q$. This
represents a structural impossibility result \emph{within the
additive cost-function class~\eqref{eq:phi_compact}}: a universal
exponent cannot exist in a biologically complete transport network
whose wall cost scales sub-linearly ($p < 1$).
\end{corollary}

\begin{corollary}[Independence from flow asymmetry]
\label{cor:asym}
For an asymmetric bifurcation where a parent vessel of flow $Q$ splits into two
daughters carrying fractions $fQ$ and $(1-f)Q$ with $f \in (0,1)$, the local
branching exponent $\alpha^*(Q,f)$ defined by $r^*(Q)^{\alpha} =
r^*(fQ)^{\alpha} + r^*((1-f)Q)^{\alpha}$ obeys the same strict boundaries:
$(5+p)/2 < \alpha^*(Q,f) < 3$.
\end{corollary}

\begin{proof}
Since $r^*$ is strictly increasing in $Q$ (Theorem~\ref{thm:unique}), define
\[
  x \;=\; \frac{r^*(fQ)}{r^*(Q)}\in(0,1), \qquad
  y \;=\; \frac{r^*((1-f)Q)}{r^*(Q)}\in(0,1).
\]
The branching equation $r^*(Q)^\alpha = r^*(fQ)^\alpha + r^*((1-f)Q)^\alpha$ is
equivalent to
\[
  G(\alpha) \;\equiv\; x^\alpha + y^\alpha \;=\; 1.
\]
Since $x,y\in(0,1)$, we have $G'(\alpha) = x^\alpha\ln x + y^\alpha\ln y < 0$
(both logarithms are strictly negative), so $G$ is \emph{strictly decreasing}.
Moreover $G(0)=2>1$ and $G(\alpha)\to0<1$ as $\alpha\to+\infty$.
By the Intermediate Value Theorem, $G(\alpha^*)=1$ has a unique solution
$\alpha^*(Q,f)>0$.

In the Murray limit ($C\to0$), $r^*\propto Q^{1/3}$, so $x=f^{1/3}$,
$y=(1-f)^{1/3}$, and $G(3)=f+(1-f)=1$; hence $\alpha^*=3$.
In the wall-dominated limit ($B\to0$), $r^*\propto Q^{2/(5+p)}$, so
$x=f^{2/(5+p)}$, $y=(1-f)^{2/(5+p)}$, and $G((5+p)/2)=1$; hence
$\alpha^*=(5+p)/2$.

For the intermediate regime $B,C>0$, $r^*(Q)$ is not a power law
(Corollary~\ref{cor:murray_uniqueness}); we evaluate the bounds explicitly.

\begin{description}
\item[\normalfont\emph{Upper bound} $G(3)<1$:\;]
Evaluate $4A = 2Br^6+(1+p)Cr^{5+p}$ at the test radius $f^{1/3}r^*(Q)$:
\[
  2B f^2 r^*(Q)^6 \;+\; (1+p)C\, f^{(5+p)/3} r^*(Q)^{5+p}.
\]
Since $p<1$, $(5+p)/3 < 2$, so $f^{(5+p)/3} > f^2$ for $f\in(0,1)$.
The expression strictly exceeds $f^2[2Br^*(Q)^6+(1+p)Cr^*(Q)^{5+p}]=4A(fQ)$.
Because the right-hand side of the optimality condition is strictly increasing
in $r$, this forces $r^*(fQ) < f^{1/3}r^*(Q)$, i.e.\ $x<f^{1/3}$, and
identically $y<(1-f)^{1/3}$, giving
\[
  G(3) \;=\; x^3 + y^3 \;<\; f + (1-f) \;=\; 1.
\]

\item[\normalfont\emph{Lower bound} $G\!\left(\tfrac{5+p}{2}\right)>1$:\;]
Evaluate at the test radius $f^{2/(5+p)}r^*(Q)$:
\[
  2B\, f^{12/(5+p)} r^*(Q)^6 \;+\; (1+p)C\, f^2 r^*(Q)^{5+p}.
\]
Since $p<1$, $12/(5+p)>2$, so $f^{12/(5+p)}<f^2$.
The expression is strictly less than $4A(fQ)$, forcing
$r^*(fQ) > f^{2/(5+p)}r^*(Q)$.
Setting $\beta=(5+p)/2$, this gives $x^\beta > f$ and $y^\beta > 1-f$, so
\[
  G\!\left(\tfrac{5+p}{2}\right) \;=\; x^\beta + y^\beta \;>\; 1.
\]
\end{description}

Since $G$ is strictly decreasing, the two bounds
$G(3)<1$ and $G\!\left(\tfrac{5+p}{2}\right)>1$ rigorously imply
$(5+p)/2 < \alpha^*(Q,f) < 3$ for all $f\in(0,1)$.
Non-universality is therefore an inherent property of the composite cost
function, not a geometric artifact of the symmetric split assumption.
\end{proof}

\end{document}
